\clearpage
\section{TX-side Conversion Gain --- Upper Sideband (No Spectral Flip)}
\label{sec:tx-gain-upper}

\textbf{What is conversion gain?} The TX side takes a signal in at FR1 and sends it out at a much higher FR3
frequency, amplifying it along the way. Conversion gain measures that amplification across the change of
frequency. Feed a signal into the ``TX-side FR1 IF In'' connector at, say, −40\,dBm, and measure the signal that comes out of the ``TX-side FR3 RF Out'' connector at its new frequency, say +5\,dBm. The conversion gain is simply the difference:
+5 − (−40) = 45\,dB. Because both powers are in dBm, subtracting them gives the conversion gain in dB directly.

\textbf{What's plotted.} Each of the first seven figures holds FR1 fixed at one center value (1.0 up to the
unit's 4.0\,GHz maximum, in 0.5\,GHz steps) and plots conversion gain against the full measured FR3 range for
that input.\footnote{Every point on both traces is a real measurement, never a simulation or an interpolation
between sparse samples: the FieldFox's built-in tracking generator drives FR1 while the FieldFox's own
receiver tracks the corresponding FR3 output in lock-step, at a calibrated −40\,dBm input level.} Grey is the
raw amplifier-and-mixer chain with the AM3153 filter bypassed; blue is the filter engaged. Both traces plot the
\emph{average} conversion gain across each point's own real swept window, nominally 500\,MHz wide (the same for
every point on every panel) --- mean linear power, converted back to dB, not a single CW tone: what a real
wideband signal occupying that window would experience. The gap between the two traces, at any point, is
exactly how much the filter costs in insertion loss there.

\textbf{Measurement conditions.} \lmxDriveRule{} (set from the LO frequency), with the LO leakage suppressed by the unit's stored LO-null calibration (Section~\ref{sec:lo-suppression}), exactly as in normal use. The unit was measured hot, as the worst case: no fan, LMX2820 die at \condGainMinC--\condGainMaxC\textdegree C. A cooler unit will typically do somewhat better.

\textbf{How the filter is chosen.} The blue trace's filter is re-picked fresh at every single plotted point,
using a two-step rule: first find whichever of the 48 candidate configurations (3 bands × 16 sub-states)
passes the most desired-sideband energy on its own, ignoring the image; then, restricted to candidates within
3\,dB of that best desired power, choose whichever rejects the image the most. Consecutive points that land on
the same filter form a \emph{region}; the dashed vertical lines on each panel mark every boundary between
regions, and the blue hex code running up from the x-axis inside each region (e.g.\ \texttt{0x3F}) is the
actual \texttt{filter\_bits} value used for every point in it --- so the exact filter behind any measurement is
always readable directly off the figure.

\textbf{Window width.} Every point averages over its own real swept window, and that window is always exactly
500\,MHz wide: a fixed filter setting's measured response across that full window is exactly what a real
wideband signal at that setting would experience.\footnote{The one exception worth noting is Figure
\ref{fig:1-7} (FR1 = 4.0\,GHz): its window is 3.75--4.25\,GHz at every point, and 250\,MHz of that exceeds the
unit's 0.7--4.0\,GHz IF design spec --- a real, measured operating point, called out directly on that panel
and outlined in red on the two heatmaps (Figures \ref{fig:1-8} and \ref{fig:1-9}).} Those two heatmaps repeat
the bypassed and optimal data across the full measured (FR1, FR3) grid, for a single-glance view instead of the seven separate slices above. FR3 values above the unit's rated 24\,GHz are boxed in red on both heatmaps
too: measured and shown, but out of design spec.

\figpair{images/sec1_fig1.pdf}{FR1 = 1.0\,GHz.}{fig:1-1}%
        {images/sec1_fig2.pdf}{FR1 = 1.5\,GHz.}{fig:1-2}
\figpair{images/sec1_fig3.pdf}{FR1 = 2.0\,GHz.}{fig:1-3}%
        {images/sec1_fig4.pdf}{FR1 = 2.5\,GHz.}{fig:1-4}
\figpair{images/sec1_fig5.pdf}{FR1 = 3.0\,GHz.}{fig:1-5}%
        {images/sec1_fig6.pdf}{FR1 = 3.5\,GHz.}{fig:1-6}
\figonesmall{images/sec1_fig7.pdf}{FR1 = 4.0\,GHz. 250\,MHz of this window falls outside the unit's IF design spec (shown in red).}{fig:1-7}
\figrow{images/sec1_fig8.pdf}{Conversion gain, upper sideband: full measured (FR1, FR3) grid, filter bypassed.
Zoom in to read each cell's exact values, printed as (FR1, FR3, conversion gain), with FR1 and FR3 in GHz and
the conversion gain in dB.}{fig:1-8}
\figrow{images/sec1_fig9.pdf}{Conversion gain, upper sideband: full measured (FR1, FR3) grid, optimal filter.
Zoom in to read each cell's exact values, printed as (FR1, FR3, conversion gain), with FR1 and FR3 in GHz and
the conversion gain in dB.}{fig:1-9}
